\section{Robustness and Scope}
\label{sec:robustness}

The switching result is not intended to be universal. Its content is that a cross-instrument sign reversal exists under economically interpretable conditions and, in the canonical Cournot region, is supported by a true global policy SPNE. This section asks whether the local sign mechanism survives alternative product-market competition and incomplete local ownership. These exercises are deliberately narrower than the baseline global-equilibrium certificate.

\subsection{Differentiated Bertrand competition}

Retain the quadratic utility in \eqref{eq:grossutility} but let firms choose prices in the final stage. The corresponding direct demand is
\begin{equation}
q_i
=\frac{a(1-\theta)-p_i+\theta p_j}{1-\theta^2},
\qquad 0\leq\theta<1.
\label{eq:bertrand-demand}
\end{equation}
The price first-order condition implies
\begin{equation}
p_i-c_i=(1-\theta^2)q_i,
\label{eq:bertrand-margin}
\end{equation}
so operating profit is $(1-\theta^2)q_i^2$. Backward induction through the investment stage gives
\begin{equation}
x_i^B=\frac{\eta(\theta)q_i}{k_x},
\qquad
g_i^B=\frac{\eta(\theta)\mu q_i+s_i}{k_g},
\end{equation}
where
\begin{equation}
\eta(\theta)=\frac{2(2-\theta^2)}{4-\theta^2}.
\end{equation}
The reduced interior government problem remains linear-quadratic.

Using the same primitive parameter vector as the Cournot witness, the Bertrand cross-infrastructure response is negative at $\theta=.3$ and positive at $\theta=.5$:
\[
\left.\frac{\partial h_A^{BR}}{\partial s_B}\right|_{\theta=0.3}
\simeq-0.00169,
\qquad
\left.\frac{\partial h_A^{BR}}{\partial s_B}\right|_{\theta=0.5}
\simeq+0.00401.
\]
The crossing occurs at approximately
\[
\theta_B^*\simeq0.397.
\]
The symmetric policy candidates are positive and locally strictly concave on the reported interval; for example, $(s,h)\simeq(3.496,1.019)$ at $\theta=.3$ and $(2.315,0.209)$ at $\theta=.5$.

The nonnegativity constraint becomes relevant not far above this range: the interior infrastructure solution reaches $h=0$ at approximately $\theta=.545$. The interior IFT must therefore not be extrapolated beyond its active-policy region. This exercise establishes that the sign-reversal mechanism is not specific to quantity competition on an open interior neighborhood; it does not establish a global Bertrand-SPNE theorem for all rivalry levels.

\subsection{Partial local capture of firm surplus}

The baseline assumes that local producer surplus is fully valued by the local government. To relax this assumption, let $\omega\in(0,1]$ be the local capture rate and define
\begin{equation}
W_i^{\omega}
=\frac12CS+\omega\pi_i-s_i g_i
-\frac{\kappa}{2}h_i^2
-\frac d2(E_i-\bar E_i)^2.
\label{eq:omega-welfare}
\end{equation}
Since $\pi_i=PS_i+s_i g_i$ on the active branch, this can equivalently be written as
\[
W_i^{\omega}
=\frac12CS+\omega PS_i+(\omega-1)s_i g_i
-\frac{\kappa}{2}h_i^2
-\frac d2(E_i-\bar E_i)^2.
\]
When $\omega=1$, the subsidy is an intrajurisdictional transfer and the baseline objective is recovered. When $\omega<1$, some firm surplus accrues outside the jurisdiction and the subsidy outlay is not fully offset by locally captured profits.

At $\omega=.9$ and the canonical primitive vector, the interior cross-infrastructure response crosses near
\[
\theta_{\omega=.9}^*\simeq0.511.
\]
The associated symmetric candidates remain positive and locally strictly concave over the reported neighborhood; for example, $(s,h)\simeq(2.385,0.714)$ at $\theta=.5$ and $(1.969,0.559)$ at $\theta=.7$. Thus the local sign reversal survives modest ownership leakage. The result is not claimed for arbitrary $\omega$, nor is this subsection a global-boundary theorem.

\subsection{Small asymmetries}

The baseline is symmetric because the paper's object is a strategic-response threshold rather than cross-sectional heterogeneity. The canonical global-SPNE gap, the relevant policy inequalities, the government Jacobian, and the switching root are all strict. Consequently sufficiently small perturbations of jurisdiction-specific primitives preserve the local qualitative configuration by continuity. We use this only as an open-set argument; we do not add a separate asymmetric closed-form model.

\subsection{Scope restrictions}

Several extensions remain outside the theory. Plant location is fixed, so the paper does not model bidding for a discrete plant location or multinational entry. Green public procurement is not a third baseline instrument. There is no hard public budget constraint mechanically linking $s_i$ and $h_i$, and there is no exogenous firm-level constraint $x_i+g_i=\bar K$. Dynamics, learning, financing constraints, heterogeneous firms, and multi-plant capital allocation are also outside the baseline.

These exclusions are substantive. The paper's contribution is not that every realistic margin can be placed in one model. The matched canonical comparison shows that adding conventional competitive investment can create a sign reversal that is absent under the same remaining primitives when that margin is removed. It does not imply that conventional investment is necessary for every possible reversal: the repository retains a no-$x$ counterexample with different environmental primitives in which a global-equilibrium sign reversal occurs.